\section{Engel Theorem the second type}

In fact we can prove a stronger version.

\begin{definition}
    A subset $S\subseteq L$ is called a \emph{Lie subset} if
    \[
        [a,b]\in S
        \qquad
        \text{for all } a,b\in S.
    \]
    Note that a Lie subset need not be a subspace, and hence need not be a Lie subalgebra.
\end{definition}

\begin{theorem}
    Let $A$ be a finite-dimensional associative algebra over $F$, and let
    \[
        L\subseteq A^{(-)}
    \]
    be a Lie subalgebra such that $A$ is generated by $L$ as an associative algebra. Suppose that
    \[
        L=\mathrm{span}(S)
    \]
    for some Lie subset $S\subseteq L$, and that every element of $S$ is nilpotent in $A$. Then $A$ is a nilpotent associative algebra.
\end{theorem}

\begin{remark}
    In the finite-dimensional situation, the existence of a uniform exponent $n$ is automatic once every element of $S$ is nilpotent in $A$. Indeed, let
    \[
        m=\dim_F A,
    \]
    and let $a\in A$ be nilpotent. Consider the left multiplication operator
    \[
        L_a:A\to A,\qquad x\mapsto ax.
    \]
    Since $a$ is nilpotent, $L_a$ is a nilpotent linear transformation on the $m$-dimensional vector space $A$. Hence
    \[
        L_a^m=0.
    \]
    But
    \[
        L_a^m=L_{a^m},
    \]
    so $a^m=0$. Therefore every nilpotent element of $A$ has nilpotency index at most $\dim_F A$. In particular, if every element of $S$ is nilpotent in $A$, then one may take
    \[
        n=\dim_F A
    \]
    and obtain
    \[
        a^n=0
        \qquad
        \text{for all } a\in S.
    \]
    Thus, in the theorem above, the hypothesis that every element of $S$ is nilpotent in $A$ already implies the existence of a common exponent $n$.
\end{remark}

\begin{theorem}[Zorn's Lemma]
    Let $(P,\leq)$ be a partially ordered set. Assume that every chain in $P$ has an upper bound in $P$. Then $P$ contains a maximal element.
\end{theorem}

\begin{proof}
    We prove the result assuming the Well-Ordering Theorem.

    Suppose, for contradiction, that $P$ has no maximal element. Then for every $x\in P$ there exists some $y\in P$ such that
    \[
        x<y.
    \]
    By the Well-Ordering Theorem, we may choose a well-ordering $\preceq$ on $P$. For each $x\in P$, let $f(x)$ be the $\preceq$-least element of the set
    \[
        \{y\in P: x<y\}.
    \]
    Then
    \[
        x<f(x)\qquad \text{for all }x\in P.
    \]

    We now define, by transfinite recursion, a family $(x_\alpha)$ indexed by ordinals. Suppose that for all $\beta<\alpha$ the elements $x_\beta$ have been defined and form a chain. By hypothesis, the chain
    \[
        C_\alpha:=\{x_\beta:\beta<\alpha\}
    \]
    has an upper bound $u_\alpha\in P$. Define
    \[
        x_\alpha:=f(u_\alpha).
    \]
    Since $u_\alpha$ is an upper bound of $C_\alpha$, for every $\beta<\alpha$ we have
    \[
        x_\beta\leq u_\alpha < x_\alpha.
    \]
    Hence $(x_\alpha)$ is a strictly increasing family:
    \[
        \beta<\alpha \implies x_\beta<x_\alpha.
    \]

    Therefore the map
    \[
        \alpha\longmapsto x_\alpha
    \]
    is injective from the class of all ordinals into the set $P$. This is impossible, since there is no injection from the class of all ordinals into a set.

    This contradiction shows that $P$ must contain a maximal element.
\end{proof}

\begin{lemma}
    Let
    \[
        \mathcal P
        =
        \left\{
        T\subseteq S \;\middle|\;
        T \text{ is a Lie subset and } a^n=0 \text{ for all } a\in T
        \right\},
    \]
    partially ordered by inclusion. Then $\mathcal P$ has a maximal element.
\end{lemma}

\begin{proof}
    We apply Zorn's lemma. Let $\{T_i\}_{i\in I}$ be a chain in $\mathcal P$, and set
    \[
        T=\bigcup_{i\in I}T_i.
    \]
    We claim that $T\in\mathcal P$.

    First, $T$ is a Lie subset. Indeed, if $a,b\in T$, then $a\in T_i$ and $b\in T_j$ for some $i,j\in I$. Since $\{T_i\}_{i\in I}$ is a chain, one of $T_i,T_j$ is contained in the other; say $T_i\subseteq T_j$. Then $a,b\in T_j$, and since $T_j$ is a Lie subset,
    \[
        [a,b]\in T_j\subseteq T.
    \]

    Second, if $a\in T$, then $a\in T_i$ for some $i\in I$, so $a^n=0$ because $T_i\in\mathcal P$.

    Thus $T\in\mathcal P$, and clearly $T$ is an upper bound of the chain. Therefore every chain in $\mathcal P$ has an upper bound. By Zorn's lemma, $\mathcal P$ has a maximal element.
\end{proof}

\begin{proof}[Proof of the theorem]
    Consider the collection
    \[
        \mathcal P=
        \left\{
        T\subseteq S
        \;\middle|\;
        T \text{ is a Lie subset of }L
        \text{ and }
        \langle T\rangle \text{ is a nilpotent associative subalgebra of }A
        \right\},
    \]
    partially ordered by inclusion. Here $\langle T\rangle$ denotes the associative subalgebra of $A$
    generated by $T$.

    We first show that every chain in $\mathcal P$ has an upper bound. Let $\{T_i\}_{i\in I}$ be a chain
    in $\mathcal P$, and set
    \[
        T=\bigcup_{i\in I}T_i.
    \]
    Since the $T_i$ form a chain, $T$ is again a Lie subset of $L$. Let $N=\dim_F A$. Since each
    $\langle T_i\rangle$ is a nilpotent subalgebra of the finite-dimensional algebra $A$, we have
    \[
        \langle T_i\rangle^N=0
        \qquad
        \text{for all }i\in I.
    \]
    Now take any $x_1,\dots,x_N\in \langle T\rangle$. Each $x_j$ belongs to $\langle T_{i_j}\rangle$
    for some $i_j\in I$. Since the family $\{T_i\}$ is totally ordered, there exists $k\in I$ such that
    \[
        T_{i_1}\cup\cdots\cup T_{i_N}\subseteq T_k.
    \]
    Hence
    \[
        x_1,\dots,x_N\in \langle T_k\rangle,
    \]
    and therefore
    \[
        x_1\cdots x_N=0.
    \]
    Thus $\langle T\rangle^N=0$, so $T\in\mathcal P$. By Zorn's lemma, $\mathcal P$ has a maximal
    element, say $S_1$.

    Set
    \[
        B=\langle S_1\rangle.
    \]
    Then $B$ is nilpotent, so there exists $t\ge 1$ such that
    \[
        B^t=0.
    \]

    We claim that if $S_1\ne S$, then there exists $s\in S\setminus S_1$ such that
    \[
        [S_1,s]\subseteq S_1.
    \]
    Assume the contrary. Choose $s_1\in S\setminus S_1$. By assumption, there exists $x_1\in S_1$ such that
    \[
        s_2=[x_1,s_1]\notin S_1.
    \]
    Since $S$ is a Lie subset, $s_2\in S$. Repeating this argument, we construct inductively elements
    \[
        x_1,x_2,\dots \in S_1,
        \qquad
        s_1,s_2,\dots \in S\setminus S_1
    \]
    such that
    \[
        s_{r+1}=[x_r,s_r]
        \qquad
        (r\ge 1).
    \]
    Thus
    \[
        s_{r+1}=\operatorname{ad}(x_r)\cdots \operatorname{ad}(x_1)(s_1).
    \]

    We now show that this process cannot continue for $2t$ steps. Indeed, for any $y_1,\dots,y_r\in B$,
    the iterated commutator
    \[
        \operatorname{ad}(y_r)\cdots \operatorname{ad}(y_1)(s_1)
    \]
    is a sum of terms of the form
    \[
        \pm u\, s_1\, v,
    \]
    where $u,v\in B$ and the total number of factors from $B$ occurring in $u$ and $v$ is $r$.
    Therefore, if $r\ge 2t$, then in each such term either $u\in B^t$ or $v\in B^t$, so the term is $0$.
    Hence
    \[
        \operatorname{ad}(y_r)\cdots \operatorname{ad}(y_1)(s_1)=0
        \qquad
        \text{for all } r\ge 2t.
    \]
    Applying this with $y_i=x_i\in S_1\subseteq B$, we obtain
    \[
        s_{2t+1}=0,
    \]
    which is impossible since $0\in S_1$. This contradiction proves the claim.

    Choose $s\in S\setminus S_1$ such that
    \[
        [S_1,s]\subseteq S_1.
    \]
    Let
    \[
        S_2=S_1\cup\{s\}.
    \]
    Then $S_2$ is a Lie subset of $L$: indeed, brackets of elements of $S_1$ remain in $S_1$,
    \([S_1,s]\subseteq S_1\), \([s,S_1]\subseteq S_1\) by antisymmetry, and \([s,s]=0\in S_1\).

    We next show that $\langle S_2\rangle$ is nilpotent. Since $[S_1,s]\subseteq S_1$, we have
    \[
        [B,s]\subseteq B.
    \]
    Hence for every $b\in B$,
    \[
        bs=sb+[b,s],
        \qquad
        [b,s]\in B.
    \]
    By repeated use of this identity, every product in the algebra generated by $B$ and $s$ can be written
    as a linear combination of terms of the form
    \[
        s^r b,
        \qquad
        b\in B,
    \]
    with $r\ge 0$. Since $s^m=0$, we may assume $0\le r\le m-1$.

    Now consider a product of length $mt$ in $\langle S_2\rangle$. After rewriting it as above, each term
    has the form $s^r b$ with $0\le r\le m-1$. Since the original length is $mt$, such a term must contain
    at least
    \[
        mt-r\ge mt-(m-1)\ge t
    \]
    factors coming from $B$. Therefore $b\in B^t=0$, and every such term is zero. It follows that
    \[
        \langle S_2\rangle^{\,mt}=0.
    \]
    Thus $\langle S_2\rangle$ is nilpotent.

    But $S_1\subsetneq S_2$ and $S_2\in\mathcal P$, contradicting the maximality of $S_1$. Hence
    $S_1=S$. Therefore
    \[
        A=\langle L\rangle=\langle S\rangle=\langle S_1\rangle
    \]
    is nilpotent.
\end{proof}

